\section{General equation}

Let 
\be
u = (u_x, u_y).
\ee
We denote the equation in dimension 2
\be
T\frac{\partial u}{\partial t} +
A_0\frac{\partial u}{\partial x} +
A_1\frac{\partial u}{\partial y} +
B_0\frac{\partial^2 u}{\partial x^2} +
B_1\frac{\partial^2 u}{\partial y^2} +
K u = f,
\label{general_equation}
\ee
We define
\be
A = (A_0, A_1), \quad 
B = (B_0, B_1),
\ee
and
\be
v = A \odot u = (A_0 u_x, A_1 u_y),\quad
w = B \odot u = (B_0 u_x, B_1 u_y), 
\ee
Then, we simplfiy \eqref{general_equation} accordingly
\be
T\frac{\partial u}{\partial t} +
\nabla \cdot v +
\nabla^2 w +
K u = f.
\label{simplified_equation}
\ee
We integrate over a given cell $j$. We denote the edges of this cell $\ell \in j$.
\be
\int_j
T\frac{\partial u}{\partial t} +
\nabla \cdot v +
\nabla^2 w +
K u = \int_j f,
\label{equation_integrated}
\ee
and proceed term by term.
We first recall the Divergence Theorem.
\be
\int_\Omega \nabla \cdot u = \int_{\partial \Omega} n \cdot u,
\ee
where n is the outward normal of the domain.
\subsection{Taylor expansion at $\ell$}
Given a field $u$, a cell $j$ and, one of its face $\ell$ and the neighbouring cell $k$ sharing the edge $\ell$ with $j$, we approximate $u_\ell$ using taylor expansion given a 2D direction $\delta$.
\be
u_{\ell + \delta} = u_\ell + \delta \nabla u_\ell + \mathcal{O}(|\delta|^2)
\ee
We define the distance between $j$ and $k$ barycenters $d_{jk} = d_{j \ell} + d_{\ell k}$, which is the sum of the distance from $j$ to $\ell$ and $\ell$ to $k$ since $j$, $k$ and $\ell$ are aligned. Since $\ell$ is at even distance from both $j$ and $k$, $d_{j \ell} = d_{\ell k}$. 

If we substitue $\delta$ with  $d_{\ell j} n_{\ell j}$ and $d_{\ell j} n_{\ell k}$ 
\be
\begin{split}
  u_{j} = u_\ell + ( n_{\ell j} d_{\ell j} ) \cdot (\nabla u)_\ell + \mathcal{O}(d_{\ell j}^2), \\
u_{k} = u_\ell - ( n_{\ell j} d_{\ell j} ) \cdot (\nabla u)_\ell + \mathcal{O}(d_{\ell j}^2).
\label{taylor_serie_at_ell}
\end{split}
\ee

\subsection{Terms that do not involve space derivatives}

We can integrate in a straight forward way the terms $f$, $Ku$ and $\dis T\frac{\partial u}{\partial t}$ from \eqref{simplified_equation}
\be
\begin{split}
\int_j f = V_j f_j \\
\int_j Ku = V_j K_j u_j \\
\int_j  T\frac{\partial u}{\partial t}= V_j T\frac{\partial u_j}{\partial t} \\
  \label{integrate_term_no_space_der}
\end{split}
\ee

\subsection{Advection term}

We then integrate the term $\nabla \cdot v$
\be
\int_j \nabla \cdot v = \int_{\partial j} n \cdot v = \sum_{\ell \in j} \int_\ell n_{\ell j} \cdot v \approx \sum_{\ell \in j} |\ell| n_{\ell j} \cdot v_\ell,
\label{integrate_nabla_v}
\ee
where $n_{\ell j}$ denotes the unit normal at $\ell$ pointing outward from cell $j$ and $|\ell|$ denotes the length of the edge $\ell$.
Summing both equations of \eqref{taylor_serie_at_ell} (substituting $u = w$) yields
\be
v_{k} + v_{j} = 2v_\ell + \mathcal{O}(d_{\ell j}^2).
\ee
And solving for $v_\ell$
\be
v_\ell = \frac{v_{k} + v_{j}}{2} + \mathcal{O}(d_{\ell j}^2).
\ee
Therefore, the equation \eqref{integrate_nabla_v} becomes
\be
\int_j \nabla \cdot v \approx \sum_{\ell \in j} |\ell| n_{\ell j} \cdot \frac{v_k + v_j}{2},
\label{integrate_nabla_v_result}
\ee

\subsection{Diffusion term}

Let us integrate the term $\nabla^2 w$. We recall that $\nabla^2 w = \nabla \cdot (\nabla w)$ and use divergence theorem to simplify the integral.
\be
\int_j \nabla^2 w = \int_{\partial j} n \cdot (\nabla w) = \sum_{\ell \in j} \int_\ell n_{\ell j} \cdot (\nabla w) \approx \sum_{\ell \in j} |\ell| n_{\ell j} \cdot (\nabla w)_\ell.
\label{integrate_nabla_square_w}
\ee
We now need to approximate $(\nabla w)_\ell$. By susbtracting the second equation of \eqref{taylor_serie_at_ell} from the first one (after substituting $u = w$) , we have
\be
w_k - w_j = 2( n_{\ell j} d_{\ell j} ) \cdot (\nabla w)_\ell + \mathcal{O}(d_{\ell j}^2).
\ee
We recall that $2d_{\ell j} = d_{j k}$ and solve for $n_{\ell j} \cdot (\nabla w)_\ell$
\be
n_{\ell j} \cdot (\nabla w)_\ell = \frac{w_k - w_j}{d_{j k}}  + \mathcal{O}(d_{\ell j}^2).
\ee
Therefore, equation \eqref{integrate_nabla_square_w} naturally becomes
\be
\dis \int_j \nabla^2 w \approx \sum_{\ell \in j} |\ell| \frac{w_k - w_j}{d_{j k}}.
\label{integrate_nabla_square_w_result}
\ee

\subsection{Time integration}
Merging \eqref{integrate_term_no_space_der}, \eqref{integrate_nabla_v_result} and \eqref{integrate_nabla_square_w_result} into \eqref{equation_integrated} yields

\be
V_j T\frac{\partial u_j}{\partial t} +
\sum_{\ell \in j} \left( |\ell| n_{\ell j} \cdot  \left(\frac{v_k + v_j}{2}\right)\right) +
\sum_{\ell \in j} |\ell| \left(\frac{w_k - w_j}{d_{j k}}\right)+
V_j K u = V_j f_j.
\label{space_integrated}
\ee
Applying a time integration on \eqref{space_integrated} while merging every known fields together ($f_j^n$ and $u_j^n$) as well as the unkown gives
\be
V_j \frac{T}{\Delta t} u_j^{n+1}  +
\sum_{\ell \in j} |\ell| \left(
 n_{\ell j} \cdot \left(\frac{v_k^{n+1} + v_j^{n+1}}{2}\right) +
\frac{w_k^{n+1} - w_j^{n+1}}{d_{j k}}
\right)+
V_j K u_j^{n+1} = V_j f_j^{n+1}
+ V_j \frac{T}{\Delta t} u_j^n
.
\label{time_integrated}
\ee
We aim to express this equation per-dimension. We will denote $d$ the index of the dimension (0 for $x$ and 1 for $y$). Therefore equation \eqref{time_integrated} becomes
\be
V_j \frac{T_d}{\Delta t} u_{jd}^{n+1}  +
\sum_{\ell \in j} |\ell| \left(
 A_d \, n_{\ell j d} \left(\frac{u_{kd}^{n+1} + u_{jd}^{n+1}}{2}\right) +
B_d \left( \frac{u_{kd}^{n+1} - u_{jd}^{n+1}}{d_{j k}} \right)
\right)+
V_j K u_{jd}^{n+1} = V_j f_{jd}^{n+1}
+ V_j \frac{T_d}{\Delta t} u_{jd}^n
.
\ee
